\chapter{Exact packet simulation and polynomial collapse}
\label{ch:exact-simulation}

\section{Gatewise invariance}
Let $P$ be an orthogonal projector and suppose the initial state obeys $P\ket{\psi_0}=\ket{\psi_0}$. If every gate satisfies
\begin{equation}
G_jP=PG_jP,
\label{eq:gatewise-invariance}
\end{equation}
then $G_j$ maps $P\Hh$ into itself. The equation is slightly more general than commutation: it requires invariance of the selected subspace, not necessarily invariance of its orthogonal complement.

\begin{theorem}[Invariant-basis simulation]
\label{thm:invariant-basis}
Let $U=G_m\cdots G_1$ and assume \cref{eq:gatewise-invariance} for every gate. If $P\ket{\psi_0}=\ket{\psi_0}$, then
\[
U\ket{\psi_0}
=
(PG_mP)\cdots(PG_1P)\ket{\psi_0}.
\]
In any basis of $P\Hh$, the restricted circuit gives exactly the same amplitudes after embedding.
\end{theorem}

\begin{proof}
Inductively, if $P\ket{\psi_{j-1}}=\ket{\psi_{j-1}}$, then
\[
G_j\ket{\psi_{j-1}}
=
G_jP\ket{\psi_{j-1}}
=
PG_jP\ket{\psi_{j-1}},
\]
so the next state remains in $P\Hh$. Iterating proves the identity.
\end{proof}

\section{Polynomial packet-simulation theorem}
\begin{theorem}[Polynomial packet-simulation theorem]
\label{thm:poly-packet}
Let $\{C_x\}$ be a polynomial-time uniform family of quantum circuits. Suppose there is an efficiently computable packet representation such that:

\begin{enumerate}
\item the initial state lies in the represented packet;
\item every gate preserves the representation;
\item the representation size is at most $\poly(|x|)$;
\item each gate update is computable to polynomial precision in polynomial time;
\item the final acceptance probability is computable from the representation in polynomial time.
\end{enumerate}

Then the language decided by the circuit family belongs to $P$.
\end{theorem}

\begin{proof}
The simulator stores a polynomial-size classical description of the evolving state. It performs polynomially many polynomial-time gate updates and evaluates the final measurement to inverse-polynomial precision. Choosing an error smaller than the amplified promise gap yields a deterministic polynomial-time decision algorithm.
\end{proof}

\begin{provedbox}
The theorem is a genuine quantum-to-classical complexity collapse for the declared subclass. It does not assume that a full $2^n$ state vector is ever constructed.
\end{provedbox}

\section{Dimension-only corollary}
A convenient special case uses an explicit basis of a packet of dimension $D_n$.

\begin{corollary}[Polynomial-dimensional invariant packet]
If $D_n\le\poly(n)$, the packet basis is efficiently indexed, every gate has a polynomial-time sparse or structured action in that basis, and the measurement is efficiently evaluable, then the circuit family is classically polynomial-time simulable.
\end{corollary}

\section{A broader compressed representation}
The state description need not be a vector of packet amplitudes. Let $\mathcal R_n$ be a polynomial-size classical representation class with an embedding into the packet Hilbert space. If initialization, gate updates, and measurement stay within $\mathcal R_n$ and are efficient, the same proof applies. This observation will later combine packet labels with stabilizer, Gaussian, or tensor-network coordinates.

\begin{figure}[H]
\centering
\includegraphics[width=.92\textwidth]{compiler_pipeline.pdf}
\caption{A complete collapse certificate includes discovery, update, leakage, measurement, and complexity information.}
\end{figure}
